%%%%%%%%%%%%%%%%%%%%%%%%%%%%%%%%%%%%%%%%%%%%%%%%%%%%%%%%%%%%%%%%%%%%%%
%%  Copyright by Wenliang Du and SEED Lab contributors.             %%
%%  This work is licensed under the Creative Commons                %%
%%  Attribution-NonCommercial-ShareAlike 4.0 International License. %%
%%  To view a copy of this license, visit                           %%
%%  http://creativecommons.org/licenses/by-nc-sa/4.0/.              %%
%%%%%%%%%%%%%%%%%%%%%%%%%%%%%%%%%%%%%%%%%%%%%%%%%%%%%%%%%%%%%%%%%%%%%%

\newcommand{\commonfolder}{../../../common-files}

\input{\commonfolder/header}
\input{\commonfolder/copyright}

\lhead{\bfseries SEED Labs -- Copy Fail Attack Lab}

\begin{document}

\begin{center}
{\LARGE Copy Fail Attack Lab}
\end{center}

\seedlabcopyright{2026}


% *******************************************
% SECTION
% *******************************************
\section{Overview}

The Linux kernel provides a number of mechanisms for moving data between
user space and kernel space. These mechanisms are security critical, because
the kernel must treat user-provided memory as untrusted. A copy operation from
user space can fail after it has already copied part of the data. If kernel
code does not handle such a partial copy correctly, stale or attacker-controlled
data may be used in a privileged context.

The Copy Fail attack, associated with CVE-2026-31431, abuses this class of
mistake in the Linux kernel's AF\_ALG AEAD interface, especially the
\texttt{algif\_aead} path. Public descriptions of the vulnerability report
that an incorrect transfer between user-controlled data and kernel-managed
state can be abused to modify the page cache of files that should be read-only
to the attacker. The attack belongs to the same general family of page-cache
corruption ideas as Dirty COW and Dirty Pipe: the attacker does not write
directly to the target file through normal file permissions, but instead
reaches a privileged kernel state where cached file data can be changed.

The objective of this lab is for students to understand the core idea behind
copy-failure vulnerabilities and why partial-copy handling is important. To keep
the lab safe and reproducible, students will not run a real kernel exploit.
Instead, they will use a user-space model that simulates the important states:
user memory, a kernel buffer, a page-cache-like object, and a privileged
program that trusts the cached data. Students will trigger a simulated copy
failure, observe how stale data is reused, and then implement and evaluate
countermeasures.

This lab covers the following topics:

\begin{itemize}[noitemsep]
\item User-kernel copy operations and partial-copy failures
\item Page cache and permission checks
\item Time-of-check/time-of-use style reasoning in kernel paths
\item Vulnerability modeling using a safe user-space simulator
\item Defense mechanisms: zeroing buffers, rejecting partial copies, and
      validating copy results
\end{itemize}

\paragraph{Readings.}
Students should read public advisories and technical writeups about Copy Fail
and CVE-2026-31431 before starting the lab. Instructors should update the links
below if newer or more authoritative references become available.

\begin{itemize}
\item Copy Fail project page: \url{https://copy.fail/}
\item NVD entry for CVE-2026-31431:
      \url{https://nvd.nist.gov/vuln/detail/CVE-2026-31431}
\item CVE record for CVE-2026-31431: \url{https://www.cve.org/CVERecord?id=CVE-2026-31431}
\item Linux kernel stable patch referenced by public advisories:
      \url{https://git.kernel.org/stable/c/a664bf3d603dc3bdcf9ae47cc21e0daec706d7a5}
\item Linux kernel AF\_ALG documentation:
      \url{https://docs.kernel.org/crypto/userspace-if.html}
\item Related background: Dirty COW, CVE-2016-5195, and Dirty Pipe,
      CVE-2022-0847
\end{itemize}

\paragraph{Lab environment.}
This lab can be conducted on the SEED Ubuntu VM or on a generic Linux machine
with a C compiler and Python 3 installed. The lab does not require students to
run a vulnerable kernel. All attack activities are conducted against the
provided user-space simulator.

\paragraph{Important safety note.}
This lab is designed to teach the mechanism behind the vulnerability, not to
provide a working exploit for a real kernel. Students should not run public
proof-of-concept code against machines that they do not own or administer.


% *******************************************
% SECTION
% *******************************************
\section{Lab Environment}

The lab setup folder contains a simulator and a few small files. The simulator
models the vulnerable kernel logic using ordinary user-space files and memory
buffers. The important files are shown below:

\begin{lstlisting}
Labsetup/
|-- Makefile
|-- simulator/
|   |-- copyfail_sim.c
|   |-- exploit_skeleton.py
|   |-- fixed_copy.c
|-- files/
|   |-- passwd.template
|   |-- victim-note.txt
|-- tests/
    |-- run-basic-tests.sh
\end{lstlisting}

In the simulator, the file \texttt{files/passwd.template} plays the role of a
protected file. Students do not modify a real \texttt{/etc/passwd} file. The
simulator loads this file into a page-cache-like object and prints the cached
content after each experiment.

\paragraph{Building the simulator.}
After downloading and unzipping the lab setup file, enter the
\texttt{Labsetup} folder and run the following commands:

\begin{lstlisting}
$ make
$ ./simulator/copyfail_sim --help
\end{lstlisting}

\paragraph{Threat model.}
The attacker is an unprivileged user. The attacker cannot open the protected
file for writing, but can interact with a simulated kernel service that accepts
data from user memory. The service has a bug: when the copy from user memory
fails partway through, it still reuses the destination buffer in a later
operation.


% *******************************************
% SECTION
% *******************************************
\section{Lab Tasks}

% -------------------------------------------
% SUBSECTION
% -------------------------------------------
\subsection{Task 1: Getting Familiar with Copy Operations}

The objective of this task is to understand why copying from user memory is
different from copying from ordinary trusted memory. User memory can be invalid,
unmapped, or changed by another thread. A kernel copy routine therefore returns
a status value, and kernel code must check that value before trusting the
destination buffer.

Run the simulator in normal mode:

\begin{lstlisting}
$ ./simulator/copyfail_sim --mode normal --input hello
\end{lstlisting}

Please answer the following questions in your report:

\begin{itemize}[noitemsep]
\item What data is copied into the simulated kernel buffer?
\item Does the simulated page cache change?
\item Which return value indicates that the copy succeeded?
\end{itemize}

Next, run the simulator with a forced copy failure:

\begin{lstlisting}
$ ./simulator/copyfail_sim --mode fail --input hello --fail-after 2
\end{lstlisting}

Please describe what happens to the destination buffer. Is it empty, partially
updated, or unchanged? Explain why this behavior is dangerous if the caller
continues to use the buffer.


% -------------------------------------------
% SUBSECTION
% -------------------------------------------
\subsection{Task 2: Observing the Vulnerable Logic}

In this task, students will inspect a simplified vulnerable code path. The
following pseudocode shows the intended pattern:

\begin{lstlisting}[language=C]
char kbuf[BUF_SIZE];

ret = copy_from_user(kbuf, user_ptr, len);
if (ret != 0) {
    return -EFAULT;
}

use_kernel_buffer(kbuf, len);
\end{lstlisting}

The vulnerable simulator intentionally uses a different pattern:

\begin{lstlisting}[language=C]
static char kbuf[BUF_SIZE];

ret = simulated_copy_from_user(kbuf, user_ptr, len);
if (ret != 0) {
    log_warning("copy failed");
}

use_kernel_buffer(kbuf, len);
\end{lstlisting}

The second pattern is dangerous because the copy result is logged but not
enforced. The buffer may contain a mixture of new data and stale data from an
earlier operation.

Run the following commands:

\begin{lstlisting}
$ ./simulator/copyfail_sim --mode normal --input AAAAAAAA
$ ./simulator/copyfail_sim --mode fail --input BBBBBBBB --fail-after 3
\end{lstlisting}

Please provide screenshots of the buffer state after both commands. Explain how
the failed copy can preserve attacker-controlled data from the previous
operation.


% -------------------------------------------
% SUBSECTION
% -------------------------------------------
\subsection{Task 3: Simulating Page-Cache Corruption}

The objective of this task is to connect the copy-failure bug to page-cache
corruption. In real systems, the page cache stores file data in memory. If a
kernel bug allows an attacker to modify cached data for a read-only file, a
privileged program may later consume the modified cached data even though the
attacker did not have permission to write the file.

In the simulator, the protected file is loaded into a page-cache-like buffer.
Students will use the vulnerable service to alter this cached copy. The
following command shows the initial cache:

\begin{lstlisting}
$ ./simulator/copyfail_sim --show-cache
\end{lstlisting}

The file contains a disabled account named \texttt{student}. Your goal is to
change the cached account line so that the simulator's login check treats the
account as enabled. Use the provided skeleton script:

\begin{lstlisting}
$ cd simulator
$ python3 exploit_skeleton.py
\end{lstlisting}

The skeleton intentionally leaves several values blank. Students need to inspect
the simulator output and decide:

\begin{itemize}[noitemsep]
\item Which offset in the cached file should be targeted?
\item Which bytes should be placed in the stale kernel buffer?
\item Where should the simulated copy fail?
\end{itemize}

After completing the skeleton, run the simulator's login check:

\begin{lstlisting}
$ ./copyfail_sim --check-login student
\end{lstlisting}

Please provide screenshots showing the cache before the attack, after the
attack, and during the login check. Explain why the file permission check did
not prevent this simulated attack.


% -------------------------------------------
% SUBSECTION
% -------------------------------------------
\subsection{Task 4: Comparing Copy Fail with Dirty COW and Dirty Pipe}

The Copy Fail attack is not identical to Dirty COW or Dirty Pipe, but all three
attacks teach an important lesson: file permissions are not the only boundary
that matters. Kernel code must also protect internal cached representations of
file data.

Please compare the three vulnerabilities at a high level. Your comparison
should include:

\begin{itemize}[noitemsep]
\item What object is corrupted or raced?
\item Does the attacker write through the normal file-write interface?
\item What kind of privileged outcome can result?
\item What assumption made by the kernel or privileged program is violated?
\end{itemize}

This is a conceptual task. Students should not download or run real exploits.


% -------------------------------------------
% SUBSECTION
% -------------------------------------------
\subsection{Task 5: Implementing Countermeasures}

In this task, students will fix the vulnerable simulator. The file
\texttt{simulator/fixed\_copy.c} contains three incomplete defense strategies.

\paragraph{Defense 1: Reject partial copies.}
The safest rule is simple: if the copy operation fails, do not use the
destination buffer. Modify the simulator so that it returns an error before any
page-cache operation occurs.

\paragraph{Defense 2: Zero the destination buffer before copying.}
Before copying user data into a reusable kernel buffer, clear the whole buffer.
This defense reduces stale-data reuse, but students should evaluate whether it
is sufficient by itself.

\paragraph{Defense 3: Use an exact valid length.}
The caller should track how many bytes were actually copied and should not pass
the original requested length to later operations after a failure.

After implementing each defense, rebuild the simulator and rerun the attack
from Task 3:

\begin{lstlisting}
$ make clean
$ make
$ python3 simulator/exploit_skeleton.py
\end{lstlisting}

Please report which defenses stop the simulated attack. If a defense does not
fully stop the attack, explain why.


% -------------------------------------------
% SUBSECTION
% -------------------------------------------
\subsection{Task 6: Kernel Hardening Discussion}

The previous task fixed a user-space model. In this task, students will reason
about real kernel hardening. Please answer the following questions:

\begin{itemize}[noitemsep]
\item Why should kernel code treat a failed user-memory copy as a security
      event rather than a minor warning?
\item Why is it dangerous to reuse static or long-lived buffers for untrusted
      user input?
\item What should regression tests check after this class of bug is fixed?
\item Why are least privilege, timely patching, and reducing unused kernel
      attack surfaces important defenses even when the exact exploit is not
      known?
\end{itemize}

Instructors may optionally ask students to inspect the patched Linux kernel
commit for CVE-2026-31431 and identify which checks were added or changed. This
optional task should be limited to code reading and should not require students
to run vulnerable kernels.


% *******************************************
% SECTION
% *******************************************
\section{Guidelines}

\paragraph{How to reason about partial copies.}
When analyzing a copy operation, do not only ask whether the copy started. Ask
whether the copy completed, how many bytes were copied, and what the destination
buffer contains after failure. The destination buffer may contain old bytes,
new bytes, zeros, or a mixture.

\paragraph{How to reason about page cache.}
The page cache is an in-memory representation of file contents. If a privileged
program reads a file and the kernel serves the data from a corrupted cache, the
program may consume attacker-influenced data even if the file on disk was not
modified through an authorized write.

\paragraph{Debugging the simulator.}
The simulator prints the buffer before and after each copy operation. When your
attack does not work, compare the following values:

\begin{itemize}[noitemsep]
\item The offset selected in the cached file
\item The bytes placed in the reusable kernel buffer
\item The failure point used by the simulated copy
\item The final cached content used by the login check
\end{itemize}

\paragraph{Instructor customization.}
Instructors can change the protected file template, the target account name,
the cache offset, the buffer size, or the required byte sequence. These changes
make the lab different each semester while preserving the same learning
objective.


% *******************************************
% SECTION
% *******************************************
\section{Submission}

%%%%%%%%%%%%%%%%%%%%%%%%%%%%%%%%%%%%%%%%
\input{\commonfolder/submission}
%%%%%%%%%%%%%%%%%%%%%%%%%%%%%%%%%%%%%%%%

In addition, your report should include:

\begin{itemize}[noitemsep]
\item The completed exploit skeleton for the simulator
\item Screenshots showing the normal copy, failed copy, cache corruption, and
      fixed behavior
\item A comparison of Copy Fail, Dirty COW, and Dirty Pipe
\item A discussion of which countermeasures worked and why
\end{itemize}

\end{document}
